\begin{solapartat}
\punts{0,25} Com que la corda està pinçada pels dos extrems, llavors tindrem un node a
cada extrem:

% DUBTE: a l'original les etiquetes de les dues figures semblen intercanviades: la figura
% de l'esquerra (3 ventres) és la que correspon al tercer harmònic i la de la dreta
% (1 ventre) al fonamental, tot i que la pauta les etiqueta a l'inrevés. Es transcriu tal
% com surt a l'original.
\begin{minipage}{0.48\linewidth}
\figura[0.9]{sol_fig1.png}
\centering Harmònic fonamental
\end{minipage}\hfill
\begin{minipage}{0.48\linewidth}
\figura[0.9]{sol_fig2.png}
\centering Tercer harmònic
\end{minipage}

\punts{0,1} Per l'harmònic fonamental $\lambda = 2L = 0,65$ m.

\punts{0,1} Pel tercer harmònic $\lambda = \dfrac{2}{3}L = 0,217$ m.

\punts{0,3} La velocitat de propagació de l'ona és: $v = \lambda f = 0,65\cdot 38,89 = 25,3$ m/s

\punts{0,3} Si mantenim la tensió constant, la velocitat de propagació no canvia:
\[ v' = v = 25,3\ \mathrm{m/s.} \]

\punts{0,1} Per l'harmònic fonamental $\lambda' = 2L = 0,54$ m.

\punts{0,1} $v = \lambda' f' \Rightarrow f' = \dfrac{v}{\lambda'} = 46,8$ Hz

Si en la representació l'estudiant ha comés l'error de no tenir en compte que la corda
està pinçada pels dos extrems i, per tant, suposa que en un extrem tenim un node però que
a l'altre extrem tenim un ventre, llavors en la resta de la resolució s'hauria d'indicar,
per exemple, que l'harmònic fonamental és $\lambda = 4L$. És a dir, que la resolució del
problema ha de ser consistent amb la representació gràfica del primer i tercer harmònic.
Cal tenir present que només es pot penalitzar una vegada un error, per tant, si la
resolució posterior és consistent amb l'error inicial, aquesta s'ha de considerar com a
correcta.
\end{solapartat}

\begin{solapartat}
\punts{0,3} $\beta = 10\log\dfrac{I}{I_0} \Rightarrow \dfrac{\beta}{10} = \log\dfrac{I}{I_0} \Rightarrow I = I_0 10^{\beta/10} = 10^{-9}\ \mathrm{W\,m^{-2}}$

\punts{0,3} $I = \dfrac{\text{Potència}}{A} \Rightarrow \text{Potència} = I\cdot A = 4\pi r^2 I = 6,65\times 10^{-6}$ W

\punts{0,2} $P' = \dfrac{\text{Potència}}{2} = 3,32\times 10^{-6}$ W

\[ I' = \frac{P'}{A} = \frac{I'}{2} = 5,00\times 10^{-10}\ \mathrm{W\,m^{-2}} \]

\punts{0,45} $\beta = 10\log\dfrac{I'}{I_0} = 10\log\dfrac{5,00\times 10^{-10}}{10^{-12}} = 27,0$ dB

\destaca{Alternativament}

\punts{0,45} $\beta = 10\log\dfrac{I'}{I_0} = 10\log\dfrac{I/2}{I_0} = 10\log\dfrac{I}{I_0} - 10\log 2 = 30 - 3,0 = 27,0$ dB
\end{solapartat}
